\documentclass{article}
\usepackage[utf8]{inputenc}
\usepackage[T1]{fontenc}
\usepackage{geometry}
\usepackage{graphicx}
\usepackage{float}
\usepackage{hyperref}
\usepackage[english,indonesian]{babel}

\title{Optimasi Pemrosesan Template LaTeX Menggunakan Pendekatan Berbasis AI}
\author{Academic Support Agent}
\date{\today}

\begin{document}
\maketitle

\begin{abstract}

\end{abstract}

\chapter{PENDAHULUAN}


\chapter{TINJAUAN PUSTAKA}


% Original content (commented out):
% \documentclass[11pt]{article}		% Précise le type de document, et la taille de la police de caractère
\usepackage[margin=2.5cm]{geometry} % Précise les marges du document
\usepackage{amsmath,amsthm,amssymb,amsfonts}	% Pour pouvoir inclure certains symboles et environnements mathématiques
\usepackage[french]{babel}	% Pour préciser la langue du document
\usepackage{graphicx}	% Pour inclure des images
\usepackage{color}	% Pour inclure du texte en couleur
\usepackage{units}	% Pour pouvoir tapper les unités correctement
\usepackage{pgf,tikz}	% Utilisation du module tikz, qui permet de tracer des belles images
\usetikzlibrary{arrows} % Quand on exporte une image GeoGebra, on a besoin de préciser cela
\usepackage{hyperref}	% Pour include des liens dans le document
\usepackage[T1]{fontenc}	% Précise la façon dont le document actuel est encodé
\usepackage[utf8]{inputenc}	% Précise comment le texte est saisi : cela permet de tapper directement les accents
\newcommand{\N}{\mathbb{N}}	% Commande personnelle, plus rapide pour tapper les ensembles
\newcommand{\Z}{\mathbb{Z}}	% Commande personnelle, plus rapide pour tapper les ensembles
\newcommand{\R}{\mathbb{R}}	% Commande personnelle, plus rapide pour tapper les ensembles
\usepackage{cprotect}	% Pour pouvoir personaliser la légende des figures
 
\begin{document}
\thispagestyle{empty}	% Pour éviter d'avoir un en-tête et un pied de page sur la page couverture
\includegraphics[width=5cm]{logo.jpg}	% Pour inclure le logo (on précise la largeur de l'image)
\vspace{4cm}	% Espacement vertical
\begin{center}	% On centre le texte
{\huge \bf Devoir 1}\\	% \huge fait que le texte est gros, \bf fait que le texte est gras
\vspace{3cm}
\large Travail présenté à Justin Bieber dans le cadre du cours \emph{MAT-2901 : Mathématiques et technologie}\\
\vspace{3cm}
Travail réalisé par Donald Trump et Donald Duck
\vfill	% On va jusqu'au bas de la page avant de mettre le texte ci-dessous
Hiver 2016
\end{center}
\newpage

\noindent \textbf{Question 1}\\
Trouvez toutes les solutions de l'équation $\sin \left(2x+\frac{\pi}{2}\right)=0$ dans $\R$.\\
% Le contenu mathématique est entre signe de $. Plutôt que d'écrire sin, on écrit \sin, ce qui met le texte droit plutôt qu'en italique. \left( et \right) permet d'ajuster la hauteur de la parenthèse selon ce qu'il y a entre elles.


\noindent \textcolor{red}{\bf Solution}\\
\noindent La fonction $\sin$ s'annule aux multiples de $\pi$. Ainsi,
\begin{align*}	% Utile pour aligner des équations. Les équations sont alignées sur le symbole &
\sin \left(2x+\frac{\pi}{2}\right)=0 &\iff 2x+\frac{\pi}{2}=k\pi,\\
		  &\iff 2x=k\pi-\frac{\pi}{2},\\
          &\iff x=\frac{k\pi}{2}-\frac{\pi}{4},
\end{align*}
où $k\in\Z$. Ainsi, l'ensemble des solutions cherchées est $\left\{\frac{k\pi}{2}-\frac{\pi}{4}:k\in\Z\right\}$.\\

\noindent \textbf{Question 2}\\
Dessinez un triangle $\Delta ABC$ tel que:
\renewcommand{\labelitemi}{$\circ$}
\begin{itemize}	% Pour une liste qui n'est pas énumérée. Si on veut des numéros, on utilise l'environnement enumerate
\item $AB=\unit[2]{cm}$;
\item $AC=\unit[\sqrt{5}]{cm}$;
\item $AC=BC$.\\
\end{itemize}
\noindent \textcolor{red}{\bf Solution}\\
En utilisant \href{http://www.texample.net/tikz/examples/}{TikZ}, on peut faire le dessin ci-dessous.
\begin{center}
\begin{tikzpicture}
\draw[thick] (0,0) -- (2,0) -- (1,2) -- (0,0);
\node[below] at (0,0) {$A$};
\node[below] at (2,0) {$B$};
\node[above] at (1,2) {$C$};
\end{tikzpicture}
\end{center}
En utilisant GeoGebra, on peut modifier les points créés aux sommets du triangle, enlever les axes, et exporter la zone de l'écran qui nous intéresse en format PDF. Lorsque c'est fait, il ne reste plus qu'à inclure l'image dans le document.
\begin{center}
\includegraphics{triangle.pdf}
\end{center}
Le triangle n'est pas de la taille voulue. On peut alors tâtonner et ajuster la largeur de l'image jusqu'à ce qu'elle soit de la taille désirée. On peut aussi retourner dans GeoGebra, et exporter au format PNG en précisant la taille de l'image en pixels.

\begin{figure}
\begin{center}
\includegraphics{triangle2.png}
\caption{Nous avons obtenu cette figure à partir de GeoGebra en exportant au format PNG en précisant la taille en pixels correspondant à la dimension réelle de l'image à l'écran.}
\end{center}
\end{figure}

\begin{figure}
\begin{center}
\begin{tikzpicture}[line cap=round,line join=round,>=triangle 45,x=1.0cm,y=1.0cm]
\clip(-0.38,-0.12) rectangle (2.26,2.31);
\draw (1,2)-- (0,0);
\draw (0,0)-- (2,0);
\draw (2,0)-- (1,2);
\begin{scriptsize}
\fill [color=black] (0,0) circle (1.0pt);
\draw[color=black] (-0.23,0.21) node {$A$};
\fill [color=black] (2,0) circle (1.0pt);
\draw[color=black] (2.15,0.21) node {$B$};
\fill [color=black] (1,2) circle (1.0pt);
\draw[color=black] (1.14,2.2) node {$C$};
\end{scriptsize}
\end{tikzpicture}
\cprotect\caption{Nous avons obtenu cette figure à partir de GeoGebra en exportant l'image en un code TikZ. Nous avons ensuite copier/coller tout ce qui est inclus à partir de \texttt{\textbackslash begin$\lbrace\text{tikzpicture}\rbrace$} jusqu'à \texttt{\textbackslash end$\lbrace\text{tikzpicture}\rbrace$}. D'autres commandes doivent parfois être copiées dans l'en-tête du document.}
\end{center}
\end{figure}
\end{document}

\end{document}